\section*{Bab \ref{ch:b3:stem-cells} --- Sel Punca, Regenerasi, dan Penuaan}

\begin{solution}{exo:b3:stem-cells:1}
\omterm{def:b3:stem-cells:stem}{Pembaruan diri}: sebuah pembelahan yang menghasilkan sedikitnya satu anakan yang
masih sebuah \omterm{def:b3:stem-cells:stem}{sel punca} --- yakni sel \emph{Lgr5} pada pangkal kriptanya. Potensi:
jangkauan nasib yang masih dapat diambil sebuah sel --- \omterm{def:b3:stem-cells:stem}{sel punca} kriptanya
multipoten, yang membuat sel penyerap, sel goblet, sel enteroendokrin, \omterm{prop:b3:stem-cells:crypt}{sel Paneth},
dan sel tuft. \omterm{def:b3:stem-cells:stem}{Ceruk}: \omterm{prop:b3:stem-cells:crypt}{sel Paneth} dan isyarat Wnt, Notch, serta EGF-nya,
yang menahan kepuncaannya di tempatnya. \omterm{def:b3:stem-cells:stem}{Sel transit-pelipat ganda}: sebuah
anakan berkomitmen yang membelah empat atau lima kali lagi di dalam kriptanya
sebelum berdiferensiasi dalam perjalanannya naik ke jonjotnya.
\end{solution}

\begin{solution}{exo:b3:stem-cells:2}
Sel sumsum tunggal menghasilkan koloni limpa yang memuat sel merah,
granulosit, dan keping darah: satu sel, beberapa garis keturunan ---
yakni multipotensi, dan tanggapan dosis yang linear menunjukkan tiap koloninya berasal dari
satu sel. \omterm{def:b3:stem-cells:stem}{Pembaruan diri} menuntut seekor mencit kedua: sel dari sebuah koloni,
yang disuntikkan ke dalam seekor penerima lain yang disinari, membuat koloni baru, sehingga
sel pendirinya telah menghasilkan anakan berkemampuan seperti dirinya.
\end{solution}

\begin{solution}{exo:b3:stem-cells:3}
\emph{Oct4}, \emph{Sox2}, \emph{Klf4}, \emph{c-Myc}. \omterm{def:b3:chromatin-epigenetics:nucleosome}{Kromatin} fibroblasnya
telah menutup gen \omterm{def:b3:stem-cells:pluripotency}{pluripotensinya} di balik \omterm{def:b3:chromatin-epigenetics:nucleosome}{heterokromatin} dan
\omterm{def:b3:chromatin-epigenetics:methylation}{metilasi DNA}, sehingga faktornya mula-mula menemukan sedikit tapak yang terjangkau lalu
harus merekrut enzim pembentuk ulang sepanjang banyak \omterm{def:b3:cell-cycle-apoptosis:cdk}{daur sel}; selnya juga harus
mematikan acaranya sendiri, lolos dari penuaan, lalu melewati sebuah
keadaan antara yang acak yang di dalamnya kebanyakan selnya macet atau mati. Hanya
sel langka yang di dalamnya tiap langkahnya berhasil yang muncul pluripoten, berminggu-minggu
kemudian.
\end{solution}

\begin{solution}{exo:b3:stem-cells:4}
Planaria menyimpan \omterm{def:b3:stem-cells:stem}{sel punca} dewasa yang pluripoten, yakni \omterm{def:b3:stem-cells:regeneration}{neoblasnya}, yang
berpindah ke lukanya lalu membangun apa pun yang hilang. Salamandernya tidak punya
kumpulan semacam itu: sel yang berdiferensiasi pada tunggulnya berdediferensiasi menjadi
\omterm{def:b3:stem-cells:stem}{pendahulu}, yang bersama \omterm{def:b3:stem-cells:stem}{sel punca} jaringan yang menetap membentuk sebuah \omterm{def:b3:stem-cells:regeneration}{blastema}
yang menjalankan ulang acara anggota badannya. Cangkok jaringan tunggal Kragl yang
ditandai menunjukkan bahwa sel \omterm{def:b3:stem-cells:regeneration}{blastemanya} tetap setia pada asalnya ---
otot membuat otot, tulang rawan tulang rawan --- sehingga \omterm{def:b3:stem-cells:regeneration}{blastemanya} sebuah
campuran \omterm{def:b3:stem-cells:stem}{pendahulu} yang terbatas garis keturunannya, bukan sebuah massa yang pluripoten.
\end{solution}

\begin{solution}{exo:b3:stem-cells:5}
$(11 - 5)/0.075 = 80$ penggandaan. Dari \qty{8}{kb}: $(8 - 5)/0.075 = 40$.
Dengan pengimbangan \qty{60}{\%}, kehilangan netonya \qty{30}{bp}: $6000/30
= 200$.
\end{solution}

\begin{solution}{exo:b3:stem-cells:6}
Empat belas pembelahan \omterm{def:b3:stem-cells:stem}{sel punca} tiap hari menghasilkan 28 anakan; 14 harus tetap
\omterm{def:b3:stem-cells:stem}{sel punca}, sehingga 14 berkomitmen --- yakni satu anakan berkomitmen tiap \omterm{def:b3:stem-cells:stem}{sel punca} tiap
hari. Masing-masing berkomitmen menjadi $2^{4} = 16$ sel: $14\times 16 = 224$ sel
jonjot tiap hari, yakni seorde dengan \num{300} yang teramati (sebuah pembelahan transit
kelima memberi \num{448}). Tingkat puncanya menyumbang pembelahan setara empat belas sel,
tingkat transitnya \num{200} sisanya.
\end{solution}

\begin{solution}{exo:b3:stem-cells:7}
$2\times 10^{11}/2^{20} = 1.9\times 10^{5}$ anakan berkomitmen tiap hari;
dengan \num{10000} \omterm{def:b3:stem-cells:stem}{sel punca}, yakni 19 tiap \omterm{def:b3:stem-cells:stem}{sel punca} tiap hari --- terhadap satu
pembelahan setahun. Cacah dua puluh pembelahannya berlaku tiap garis keturunan dari
\omterm{def:b3:stem-cells:stem}{sel punca} ke sel merah, tetapi \omterm{def:b3:stem-cells:stem}{pendahulu} antaranya bukanlah
sekali pakai: kumpulan \omterm{def:b3:stem-cells:stem}{pendahulu} multipoten dan berkomitmen menopang dirinya
berminggu-minggu lewat pembelahannya sendiri, sehingga nyaris semua pembelahannya terjadi
di tingkat \omterm{def:b3:stem-cells:stem}{pendahulunya} dan \omterm{def:b3:stem-cells:stem}{sel puncanya} jarang dipanggil.
\end{solution}

\begin{solution}{exo:b3:stem-cells:8}
$\mu(50) = 10^{-3}\mathrm{e}^{1.74} = 5.7\times 10^{-3}$; $\mu(70) =
10^{-3}\mathrm{e}^{3.48} = 0.032$; $\mu(90) = 10^{-3}\mathrm{e}^{5.22} =
0.19$ tiap tahun. Dengan $\mu_{0}/\gamma = 0.0115$: $S(70) = \exp[-0.0115
\times 31.5] = 0.70$, $S(90) = \exp[-0.0115\times 184] = 0.12$. Mediannya:
$30 + \ln(61)/0.087 = 77$. Membagi dua $\mu_{0}$: $\ln(121)/0.087 = 55$,
yakni median 85 --- satu waktu penggandaan, delapan tahun, yang diperoleh. Membagi dua $\gamma$
menjadi $0.0435$: $\ln(31)/0.0435 = 79$, yakni median 109: memperlambat pangkatnya
memberi empat kali lebih banyak daripada membagi dua tetapannya.
\end{solution}

\begin{solution}{exo:b3:stem-cells:9}
Sepertiga yang tersisa harus melipattigakan diri: $\log_{2}3 = 1.6$ pembelahan tiap sel.
\omterm{def:b3:stem-cells:regeneration}{Pertumbuhan pengimbang}, sebab lobus yang ada membesar lewat pembelahan
sel di tempatnya; lobus yang dibuang, dengan saluran, pembuluh, dan
arsitekturnya, tidak dibangun kembali --- hatinya memulihkan massa dan
fungsinya tetapi bukan bentuknya.
\end{solution}

\begin{solution}{exo:b3:stem-cells:10}
Pada tiap peristiwanya, ukuran klonanya $k$ beranjak ke $k + 1$ atau $k - 1$ dengan peluang
yang sama, sehingga harapan ukurannya tidak pernah berubah; ketika jalannya berakhir, ukurannya
$N$ (pengambilalihan) atau $0$ (kehilangan), dan harapannya $N\,P + 0 = k$
memberi $P = k/N$. Sebuah sel tunggal: $1/14$. Sebuah mutasi netral pada satu \omterm{def:b3:stem-cells:stem}{sel
punca} terpancang di dalam kriptanya dengan peluang $1/14$ lalu jika tidak hilang
dalam beberapa bulan --- hanyutannya membuang tiga belas dari empat belas mutasinya.
\end{solution}

\begin{solution}{exo:b3:stem-cells:11}
Jalannya kini berat sebelah: klonanya memperoleh sebuah tempat lebih sering daripada ia
kehilangan satu, sehingga harapan ukurannya tumbuh dan peluang pengambilalihannya
naik ke arah kepastian seiring tumbuhnya $k$ --- bagi sebuah keunggulan \qty{10}{\%} di
sebuah \omterm{def:b3:stem-cells:stem}{ceruk} beranggota empat belas, kira-kira tiga kali lipat di atas $1/14$ bagi sebuah sel tunggal,
dan lebih tinggi lagi begitu klonanya memegang beberapa tempat. Selama berpuluh tahun, \omterm{def:b3:stem-cells:stem}{sel
punca} dengan mutasi semacam itu (\emph{DNMT3A}, \emph{TET2} di sumsumnya) mengambil
alih \omterm{def:b3:stem-cells:stem}{ceruknya}, dan pada umur tujuh puluh seperlima orang punya sebagian besar
darahnya dari satu klona. Ia bukanlah \omterm{def:b3:cancer-biology:hallmarks}{kanker}: selnya masih
berdiferensiasi lalu menaati hierarkinya. Tetapi mutasi berikutnya kini punya
sebuah klona bermiliar-miliar untuk timbul di dalamnya, dan risiko leukemianya naik
sepuluh kali lipat.
\end{solution}

\begin{solution}{exo:b3:stem-cells:12}
Mediannya $= 30 + \gamma^{-1}\ln(1 + \gamma\ln 2/\mu_{0})$: bagi $\mu_{0} =
10^{-2}$, $30 + \ln 7/0.087 = 52$; $10^{-3}$: 77; $10^{-4}$: $30 +
\ln 604/0.087 = 104$. Pada bentuk Gompertz yang murni, tiap pemangkasan sepuluh kali lipatnya menambahkan
$\ln 10/\gamma = 26$ tahun yang sama. Hasilnya menyusut pada praktiknya
sebab perolehan sejarahnya datang dari menyingkirkan sebuah suku yang tak bergantung umur
--- infeksi, persalinan, kecelakaan, tetapan Makeham $A$ pada $\mu = A
+ \mu_{0}\mathrm{e}^{\gamma t}$ --- dan begitu $A$ kecil dibandingkan
eksponensialnya pada umur yang berarti, memangkasnya lebih jauh nyaris tidak mengubah
apa pun; $\mu_{0}$ yang hakikinya bergerak jauh lebih sedikit. Sebuah pemangkasan \qty{20}{\%}
pada $\gamma$ (menjadi $0.070$) memberi $30 + \ln 49/0.070 = 86$: yakni sembilan tahun,
yang diperoleh pada tiap umurnya alih-alih dengan menunda beberapa sebab kematiannya.
\end{solution}

\begin{solution}{pb:b3:stem-cells:1}
\textbf{1.} $3.5\times 10^{11}$ tiap hari, $4\times 10^{6}$ tiap detik.
\textbf{2.} $2^{20} \approx 10^{6}$ sel tiap anakan berkomitmen;
$3.5\times 10^{11}/10^{6} = 3.3\times 10^{5}$ anakan berkomitmen tiap hari.
\textbf{3.} $33$ tiap \omterm{def:b3:stem-cells:stem}{sel punca} tiap hari, \num{12000} tiap tahun.
\textbf{4.} Kumpulan \omterm{def:b3:stem-cells:stem}{pendahulunya} harus menopang dirinya: \omterm{def:b3:stem-cells:stem}{pendahulu} multipoten
dan berkomitmen membelah berminggu-minggu lalu menjaga cacahnya sendiri,
sehingga nyaris tiap pembelahan dari kedua puluhnya terjadi di tingkat yang membarui
dirinya, dan \omterm{def:b3:stem-cells:stem}{sel puncanya} jarang menyumbang sebuah anakan berkomitmen
--- yakni berorde sekali setahun.
\textbf{5.} $10 - 20\times 0.06 = \qty{8.8}{kb}$: jauh di atas
$L_{\text{crit}}$, dan sel merahnya tidak pernah membelah lagi; cadangannya
berarti hanya bagi \omterm{def:b3:stem-cells:stem}{pendahulu} yang dipakai ulang.
\textbf{6.} Kehilangan netonya $0.3\times 60 = \qty{18}{bp}$; $6000/18 = 333$
pembelahan; pada satu tiap tahun, yakni $333$ tahun --- jauh melampaui sebuah hidup, sehingga
telomernya tidak membatasi \omterm{def:b3:stem-cells:stem}{sel puncanya} sendiri; kerusakan lain yang membatasinya.
\textbf{7.} $14$ pembelahan tiap hari; $7$ menggantikan \omterm{def:b3:stem-cells:stem}{sel punca} yang hilang dan $7$
berkomitmen; $7\times 2^{4} = 112$ sel jonjot tiap hari.
\textbf{8.} $N^{2} = 196$ peristiwa tiap sel pada satu tiap hari: yakni sekitar
\qty{200}{d}, enam sampai tujuh bulan --- yakni waktu yang sepanjangnya kripta
konfetinya menjadi berwarna tunggal.
\textbf{9.} $1/14 \approx \qty{7}{\%}$; $10^{7}/14 \approx 7\times 10^{5}$
kripta.
\textbf{10.} Hanyutannya membuang tiga belas dari tiap empat belas mutasinya
dalam beberapa bulan. \omterm{def:b3:stem-cells:stem}{Pembelahan tak setangkup} seumur hidup akan menyimpan tiap \omterm{def:b3:stem-cells:stem}{sel
punca} mutan seumur hidup, sehingga mutasinya akan menimbun di tiap garis keturunannya
tanpa pernah dibuang.
\textbf{11.} Kepuncaan adalah sebuah kedudukan di dalam \omterm{def:b3:stem-cells:stem}{ceruknya} dan sehimpunan isyarat,
bukan sebuah sifat tetap sebuah sel: sel mana pun yang mencapai isyarat
Panethnya dapat mengambil perannya.
\textbf{12.} Imbaslah sebuah label berwarna banyak pada \omterm{def:b3:stem-cells:stem}{sel puncanya} pada satu saat
lalu ikutilah kriptanya. Hanyutan setangkup: kriptanya menjadi berwarna tunggal
dalam beberapa bulan. \omterm{def:b3:stem-cells:stem}{Pembelahan tak setangkup}: tiap kriptanya menyimpan semua warnanya
tanpa batas, masing-masing sebagai sebuah pita tipis naik ke jonjotnya.
\textbf{13.} $6000/60 = 100$ penggandaan.
\textbf{14.} Sekitar $75$ tersisa: selnya mencacah pembelahan, bukan waktu
penanggalan, sebab sedasawarsa di dalam pembekunya tidak berongkos apa pun.
\textbf{15.} $(8000 - 4000)/30 = 133$ tahun sesudah umur dua puluh --- yakni umur 153.
Panjang telomer leukosit rata-ratanya tidak dengan sendirinya membatasi hidupnya;
telomer terpendeknya, dan ciri khas lainnya, lebih berarti.
\textbf{16.} \omterm{def:b3:cancer-biology:telomerase}{Telomerase} menyingkirkan satu penghalang, yakni \omterm{thm:b3:stem-cells:hayflick}{penuaan replikatif},
sehingga ia ditemukan pada nyaris tiap \omterm{def:b3:cancer-biology:hallmarks}{kanker}, yang memerlukan pembelahan
tanpa batas; tetapi fibroblas yang diberi \omterm{def:b3:cancer-biology:telomerase}{telomerase} menyimpan titik periksanya,
penghambatan sentuhannya, dan p53-nya yang utuh, sehingga pembelahan tanpa batas saja tidak
menjadikannya \omterm{def:b3:cancer-biology:hallmarks}{kanker}. Perlu, tidak cukup.
\textbf{17.} $\mu(40) = 2.4\times 10^{-3}$, $\mu(60) = 0.014$, $\mu(80) =
0.077$, $\mu(100) = 0.44$ tiap tahun. $S = \exp[-0.0115(\mathrm{e}^{\gamma
(t-30)} - 1)]$: $0.98$ pada 40, $0.87$ pada 60, $0.42$ pada 80, $0.006$ pada
100.
\textbf{18.} Mediannya $30 + \ln 61/0.087 = 77$. Sepuluh persen bertahan hidup pada
$\mathrm{e}^{\gamma(t-30)} = 1 + \gamma\ln 10/\mu_{0} = 201$: $t = 30 +
5.3/0.087 = 91$.
\textbf{19.} \qty{0.67}{mm/d}; $\log_{2}1000 = 10$ penggandaan dalam 60
hari, yakni satu tiap enam hari.
\textbf{20.} Sekitar $10^{8}$ pembelahan, sehingga $0.1$ mutasi onkogenik tiap
\omterm{def:b3:stem-cells:regeneration}{regenerasi} --- yakni satu dari sepuluh anggota badan. Aksolotl nyaris tidak pernah
mengembangkan \omterm{def:b3:cancer-biology:hallmarks}{tumor}: berdarah dingin dan lambat, dengan penekanan \omterm{def:b3:cancer-biology:hallmarks}{tumor} yang tangguh, dan
sel \omterm{def:b3:stem-cells:regeneration}{blastemanya} tetap terbatas garis keturunannya lalu berdiferensiasi ulang alih-alih
terus membelah; jaringan yang beregenerasi bahkan menekan \omterm{def:b3:cancer-biology:hallmarks}{tumor}
yang diimbas. Seekor mamalia punya lebih banyak sel, suhu dan laju metabolisme
yang lebih tinggi, serta hidup yang lebih panjang yang sepanjangnya sebuah klona yang lolos dapat tumbuh.
\textbf{21.} Buanglah saraf dari tunggulnya lalu tanamkanlah sebuah manik yang melepaskan
protein calonnya, atau cangkokkanlah jaringan yang dikondisikan saraf: \omterm{def:b3:stem-cells:regeneration}{regenerasi}
yang pulih berarti ada sebuah faktor yang dapat membaur. Calon yang dikenali pada kadal air dan
aksolotl: protein landaian depannya nAG, neuregulin-1, dan FGF.
\textbf{22.} \omterm{prop:b3:stem-cells:crypt}{Pengisyaratan Wnt} tinggi pada ekornya dan rendah pada kepalanya, dan
luka pada ujung depan sebuah keping biasanya membuat sebuah kepala sebab Wnt
di sana rendah; menidakaktifkan sebuah penghambat Wnt menaikkan Wnt di mana-mana,
kedua lukanya membaca ``belakang,'' dan dua ekor tumbuh. Kepingnya mengetahui
sumbunya sendiri dari landaian sisa di dalam ototnya, dan lukanya
menegakkan kembali landaiannya relatif terhadapnya.
\textbf{23.} Sepertiga yang tersisa harus melipattigakan diri: tiap hepatositnya membelah
sekali (menjadi dua pertiga) dan separuhnya membelah lagi --- yakni rata-rata $1.6$ pembelahan.
Lobus yang selamat membengkak sampai massa aslinya; lobus yang hilang serta
saluran dan pembuluhnya tidak dibangun kembali.
\textbf{24.} \omterm{def:b3:stem-cells:regeneration}{Regenerasi} yang menggantikan jaringan yang aus dan rusak akan
memperlambat penimbunan yang diukur pangkatnya, yang menurunkan $\gamma$ lalu
mendataran kurvanya --- aksolotl menunjukkan sedikit sekali penuaan. Harganya
adalah sebuah suku tetapan yang lebih tinggi dari \omterm{def:b3:cancer-biology:hallmarks}{kanker} pada sebuah tubuh yang hangat, besar, dan
berumur panjang.
\textbf{25.} Cacah Hayflicknya sekitar $100$ penggandaan; sebuah \omterm{def:b3:stem-cells:stem}{sel punca} yang dibantu
\omterm{def:b3:cancer-biology:telomerase}{telomerase} punya sekitar $330$ pembelahan, yakni lebih dari tiga abad pada satu tiap
tahun; masa hidup mediannya $77$ tahun.
\end{solution}
